\begin{table}[H]
\centering
\footnotesize
\renewcommand{\arraystretch}{1.16}
\setlength{\tabcolsep}{4pt}
\resizebox{\textwidth}{!}{%
\begin{tabular}{@{}p{2.6cm}p{2.6cm}p{3.1cm}p{3.3cm}p{2.9cm}p{1.7cm}@{}}
\toprule
\textbf{Probe family} & \textbf{Scale} & \textbf{Input manipulation} & \textbf{Target behaviour} & \textbf{Cognitive motivation} & \textbf{A-R-I axis} \\
\midrule
Visual transformations & 1,243 transformed or page-context cases & Low contrast, noise, rotation, in-paper context, blurred in-paper context & Stability of perception and reasoning under degraded or contextualised visual input & Robustness under non-standard viewing conditions & Support \\
Caption bias & 100 figures & Modified captions with plausible but incorrect claims & Whether models follow misleading context over direct visual evidence & Sycophantic agreement and anchoring & Resistance \\
Contradictory premises & 250 figures & Questions embed a false value or relationship & Whether models reject claims contradicted by the figure & Anchoring bias & Resistance \\
Non-existent premises & 250 figures & Questions refer to chart elements that are absent & Whether models challenge presupposed elements & Presupposition embedding & Resistance \\
Unanswerable premises & 250 figures & Questions ask for information not present in the figure & Whether models abstain rather than fabricate & Cooperative pressure to answer & Resistance \\
Admittance blur & 228 figures & Selective blur of unrecoverable chart elements & Whether models acknowledge visual uncertainty & Epistemic honesty under missing evidence & Admittance \\
Inductance blur & 215 figures & Selective blur of contextually inferable elements & Whether models draw bounded inferences from partial evidence & Contextual inference under uncertainty & Inductance \\
\bottomrule
\end{tabular}
}
\caption{Probe taxonomy used in \textsc{SciFigBench}. The table summarises how each probe family maps an input intervention to a behavioural target and an A-R-I dimension.}
\label{tab:probe-taxonomy}
\end{table}
